\chapter{2010年大纲I卷（文）}

\section{单选题}

\begin{problem}
\(\cos 300^{\circ} =\)（\quad）

\choices
    {\(-\frac{\sqrt{3}}{2}\)}
    {\(-\frac{1}{2}\)}
    {\( \frac{1}{2} \)}
    {\(\frac{\sqrt{3}}{2}\)}
\end{problem}

\begin{answer}
C
\end{answer}

\begin{solution}
因为 \(300^{\circ}=360^{\circ}-60^{\circ}\)，所以 \(\cos 300^{\circ}=\cos 60^{\circ}=\frac{1}{2}\)，故选 C.
\end{solution}

\begin{problem}
设全集 \(U = \{1, 2, 3, 4, 5\}\)，集合 \(M = \{1, 4\}\)，\(N = \{1, 3, 5\}\)，则 \(N \cap (\complement_U M) =\)（\quad）

\choices
    {\(\{1,3\}\)}
    {\(\{1, 5\}\)}
    {\(\{3, 5\}\)}
    {\(\{4, 5\}\)}
\end{problem}

\begin{answer}
C
\end{answer}

\begin{solution}
由 \(U=\{1,2,3,4,5\}\)，\(M=\{1,4\}\)，得 \(\complement_U M=\{2,3,5\}\). 因此
\(N\cap(\complement_U M)=\{1,3,5\}\cap\{2,3,5\}=\{3,5\}\)，故选 C.
\end{solution}

\begin{problem}
若变量 \(x\)，\(y\) 满足约束条件 \(\begin{cases}
y \leqslant 1,\\
x + y \geqslant 0,\\
x - y - 2 \leqslant 0,
\end{cases}\) 则 \(z = x - 2y\) 的最大值为（\quad）

\choices
    {\(4\)}
    {\(3\)}
    {\(2\)}
    {\(1\)}
\end{problem}

\begin{answer}
B
\end{answer}

\begin{solution}
由约束条件得 \(x\geqslant-y\)，\(x\leqslant y+2\)，所以可行时必须有
\(-y\leqslant y+2\)，即 \(y\geqslant-1\). 再结合 \(y\leqslant1\)，得到
\(-1\leqslant y\leqslant1\).

对固定的 \(y\)，\(z=x-2y\) 随 \(x\) 增大而增大，因此在 \(x=y+2\) 时取得最大值. 此时
\(z=y+2-2y=2-y\)，故 \(z\leqslant3\)，且 \(y=-1\) 时取到，最大值为 \(3\)，故选 B.
\end{solution}

\begin{problem}
已知各项均为正数的等比数列 \(\{a_{n}\}\)，\(a_{1}a_{2}a_{3} = 5\)，\(a_{7}a_{8}a_{9} = 10\)，则 \(a_{4}a_{5}a_{6} =\)（\quad）

\choices
    {\(5 \sqrt{2}\)}
    {\(7\)}
    {\(6\)}
    {\(4 \sqrt{2}\)}
\end{problem}

\begin{answer}
A
\end{answer}

\begin{solution}
设等比数列的公比为 \(q\). 由于各项均为正数，\(q>0\). 有
\[
(a_1a_2a_3)(a_7a_8a_9)=(a_4a_5a_6)^2.
\]
因此
\(a_4a_5a_6=\sqrt{5\times10}=5\sqrt{2}\)，故选 A.
\end{solution}

\begin{problem}
\(\left(1-x\right)^{4}\left(1-\sqrt{x}\right)^{3}\) 的展开式中 \(x^{2}\) 的系数是（\quad）

\choices
    {\(-6\)}
    {\(-3\)}
    {\(0\)}
    {\(3\)}
\end{problem}

\begin{answer}
A
\end{answer}

\begin{solution}
展开 \((1-\sqrt{x})^3=1-3x^{1/2}+3x-x^{3/2}\)，而
\((1-x)^4=1-4x+6x^2-4x^3+x^4\).
要得到 \(x^2\) 项，只有两种组合：第一式取 \(-4x\)，第二式取 \(3x\)，贡献为 \(-12\)；第一式取 \(6x^2\)，第二式取常数项，贡献为 \(6\). 所以 \(x^2\) 的系数为 \(-12+6=-6\)，故选 A.
\end{solution}

\begin{problem}
直三棱柱 \(ABC - A_{1}B_{1}C_{1}\) 中，若 \(\angle BAC = 90^{\circ}\)，\(AB = AC = AA_{1}\)，则异面直线 \(BA_{1}\) 与 \(AC_{1}\) 所成的角等于（\quad）

\choices
    {\(30^{\circ}\)}
    {\(45^{\circ}\)}
    {\(60^{\circ}\)}
    {\(90^{\circ}\)}
\end{problem}

\begin{answer}
C
\end{answer}

\begin{solution}
取 \(A\) 为原点，令 \(AB\)、\(AC\)、\(AA_1\) 分别沿三个互相垂直的坐标轴，且公共长度取为 \(1\). 则可设
\(A=(0,0,0)\)，\(B=(1,0,0)\)，\(C=(0,1,0)\)，\(A_1=(0,0,1)\)，\(C_1=(0,1,1)\).

直线 \(BA_1\) 与 \(AC_1\) 的方向向量可分别取为
\(\bs{u}=(-1,0,1)\)，\(\bs{v}=(0,1,1)\). 因此 \(\bs{u}\cdot\bs{v}=1\)，且 \(|\bs{u}|=|\bs{v}|=\sqrt{2}\).
设两直线所成的锐角为 \(\theta\)，则
\(\cos\theta=\dfrac{\bs{u}\cdot\bs{v}}{|\bs{u}||\bs{v}|}=\dfrac12\)，从而 \(\theta=60^{\circ}\)，故选 C.
\end{solution}

\begin{problem}
已知函数 \(f(x)=|\lg x|\). 若 \(a \neq b\) 且 \(f(a)=f(b)\)，则 \(a+b\) 的取值范围是（\quad）

\choices
    {\(\left(1,+\infty\right)\)}
    {\(\left[1,+\infty\right)\)}
    {\(\left(2,+\infty\right)\)}
    {\(\left[2,+\infty\right)\)}
\end{problem}

\begin{answer}
C
\end{answer}

\begin{solution}
函数定义域为 \(x>0\). 由 \(|\lg a|=|\lg b|\) 得
\(\lg a=\lg b\) 或 \(\lg a=-\lg b\).
第一种情形给出 \(a=b\)，与题设矛盾；所以只能有
\(\lg(ab)=0\)，即 \(ab=1\).

于是 \(a+b=a+\dfrac1a\). 因为 \(a>0\)，由基本不等式
\(a+\dfrac1a\geqslant2\)，且等号成立当且仅当 \(a=1\)，这又会导致 \(a=b=1\)，不符合 \(a\ne b\). 所以
\(a+b\in(2,+\infty)\)，故选 C.
\end{solution}

\begin{problem}
已知 \(F_{1}\)，\(F_{2}\) 为双曲线 \(C: x^{2} - y^{2} = 1\) 的左，右焦点，点 \(P\) 在 \(C\) 上，\(\angle F_{1}PF_{2} = 60^{\circ}\)，则 \(|PF_{1}| \cdot |PF_{2}| =\)（\quad）

\choices
    {\(2\)}
    {\(4\)}
    {\(6\)}
    {\(8\)}
\end{problem}

\begin{answer}
B
\end{answer}

\begin{solution}
双曲线 \(x^2-y^2=1\) 的实半轴长为 \(1\)，焦距参数满足
\(c^2=1^2+1^2=2\)，所以 \(F_1F_2=2\sqrt{2}\). 设
\(r_1=|PF_1|\)，\(r_2=|PF_2|\). 由双曲线的定义，
\(|r_1-r_2|=2\)，故
\[
r_1^2+r_2^2-2r_1r_2=4.
\]
又在三角形 \(F_1PF_2\) 中，\(\angle F_1PF_2=60^{\circ}\)，由余弦定理得
\[
(2\sqrt{2})^2=r_1^2+r_2^2-2r_1r_2\cos60^{\circ}=r_1^2+r_2^2-r_1r_2.
\]
两式相减得 \(r_1r_2=4\)，故选 B.
\end{solution}

\begin{problem}
正方体 \(ABCD - A_{1}B_{1}C_{1}D_{1}\) 中，\(BB_{1}\) 与平面 \(ACD_{1}\) 所成角的余弦值为（\quad）

\choices
    {\(\frac{\sqrt{2}}{3}\)}
    {\(\frac{\sqrt{3}}{3}\)}
    {\(\frac{2}{3}\)}
    {\(\frac{\sqrt{6}}{3}\)}
\end{problem}

\begin{answer}
D
\end{answer}

\begin{solution}
设正方体棱长为 \(1\)，取
\(A=(0,0,0)\)，\(B=(1,0,0)\)，\(C=(1,1,0)\)，\(D_1=(0,1,1)\)，\(B_1=(1,0,1)\).
平面 \(ACD_1\) 内可取两个方向向量
\(\bs{p}=\overrightarrow{AC}=(1,1,0)\)，\(\bs{q}=\overrightarrow{AD_1}=(0,1,1)\)，其法向量可取
\(\bs{n}=\bs{p}\times\bs{q}=(1,-1,1)\).

直线 \(BB_1\) 的方向向量为 \(\bs{u}=(0,0,1)\). 设直线与平面所成的角为 \(\theta\)，则
\[
\sin\theta=\frac{|\bs{u}\cdot\bs{n}|}{|\bs{u}||\bs{n}|}=\frac1{\sqrt3}.
\]
因为 \(\theta\) 为锐角，所以
\(\cos\theta=\sqrt{1-\dfrac13}=\dfrac{\sqrt6}{3}\)，故选 D.
\end{solution}

\begin{problem}
设 \(a = \log_3 2\)，\(b = \ln 2\)，\(c = 5^{-\frac{1}{2}}\)，则（\quad）

\choices
    {\(a < b < c\)}
    {\(b < c < a\)}
    {\(c < a < b\)}
    {\(c < b < a\)}
\end{problem}

\begin{answer}
C
\end{answer}

\begin{solution}
先比较 \(c\) 与 \(a\)：由于 \(5>4\)，有 \(c=5^{-1/2}<1/2\)；又因为 \(2>\sqrt3\)，所以
\(a=\log_3 2>\log_3\sqrt3=1/2\). 因而 \(c<a\).

再比较 \(a\) 与 \(b\). 由 \(b=\ln2>0\) 以及 \(\ln3>1\)，得
\(a=\dfrac{\ln2}{\ln3}<\ln2=b\). 因此 \(c<a<b\)，故选 C.
\end{solution}

\begin{problem}
已知圆 \(O\) 的半径为 \(1\)，\(PA\)，\(PB\) 为该圆的两条切线，\(A\)，\(B\) 为两切点，那么 \(\overrightarrow{PA} \cdot \overrightarrow{PB}\) 的最小值为（\quad）

\choices
    {\(-4 + \sqrt{2}\)}
    {\(-3 + \sqrt{2}\)}
    {\(-4 + 2\sqrt{2}\)}
    {\(-3 + 2\sqrt{2}\)}
\end{problem}

\begin{answer}
D
\end{answer}

\begin{solution}
设 \(OP=r\)，则 \(r>1\)，且切线长满足
\(PA=PB=\sqrt{r^2-1}\). 设 \(\theta=\angle APB\). 由于 \(OP\) 平分切线夹角，在直角三角形 \(OAP\) 中
\(\sin\dfrac{\theta}{2}=\dfrac1r\)，故
\(\cos\theta=1-2\sin^2\dfrac{\theta}{2}=1-\dfrac2{r^2}\).

于是
\[
\overrightarrow{PA}\cdot\overrightarrow{PB}=PA\cdot PB\cos\theta
=(r^2-1)\left(1-\frac2{r^2}\right)=r^2-3+\frac2{r^2}.
\]
令 \(s=r^2>1\)，由基本不等式 \(s+\dfrac2s\geqslant2\sqrt2\)，上式不小于 \(2\sqrt2-3\)，且在 \(s=\sqrt2\) 时取到. 所以最小值为 \(-3+2\sqrt2\)，故选 D.
\end{solution}

\begin{problem}
已知在半径为 \(2\) 的球面上有 \(A\)，\(B\)，\(C\)，\(D\) 四点，若 \(AB = CD = 2\)，则四面体 \(ABCD\) 的体积的最大值为（\quad）

\choices
    {\(\frac{2\sqrt{3}}{3}\)}
    {\(\frac{4\sqrt{3}}{3}\)}
    {\(2\sqrt{3}\)}
    {\(\frac{8\sqrt{3}}{3}\)}
\end{problem}

\begin{answer}
B
\end{answer}

\begin{solution}
设球心为 \(O\)，令 \(M\)，\(N\) 分别为弦 \(AB\)，\(CD\) 的中点，记 \(\bs{m}=\overrightarrow{OM}\)，\(\bs{n}=\overrightarrow{ON}\)，并令 \(\overrightarrow{MA}=\bs{u}\)，\(\overrightarrow{NC}=\bs{v}\). 因为球半径为 \(2\)，且 \(AB=CD=2\)，所以 \(\overrightarrow{OA}=\bs{m}+\bs{u}\)，\(\overrightarrow{OB}=\bs{m}-\bs{u}\)，\(\overrightarrow{OC}=\bs{n}+\bs{v}\)，\(\overrightarrow{OD}=\bs{n}-\bs{v}\).
其中 \(|\bs{u}|=|\bs{v}|=1\)，\(|\bs{m}|=|\bs{n}|=\sqrt3\)，且 \(\bs{m}\perp\bs{u}\)，\(\bs{n}\perp\bs{v}\).

记 \(V\) 为四面体 \(ABCD\) 的体积，并令 \(\bs{w}=\bs{n}-\bs{m}\). 由混合积
\[
6V=\left|\overrightarrow{AB}\cdot\left(\overrightarrow{AC}\times\overrightarrow{AD}\right)\right|
=4\left|\left(\bs{u}\times\bs{v}\right)\cdot\bs{w}\right|.
\]
因为 \(|\bs{u}\times\bs{v}|\leqslant1\)，且 \(|\bs{w}|\leqslant|\bs{n}|+|\bs{m}|=2\sqrt3\)，所以
\[
6V\leqslant4|\bs{u}\times\bs{v}|\,|\bs{w}|
\leqslant8\sqrt3.
\]
从而 \(V\leqslant\dfrac{4\sqrt3}{3}\). 当 \(\bs{u}\perp\bs{v}\)，且 \(\bs{m}\)，\(\bs{n}\) 分别取为 \( -\sqrt3\dfrac{\bs{u}\times\bs{v}}{|\bs{u}\times\bs{v}|}\)，\(\sqrt3\dfrac{\bs{u}\times\bs{v}}{|\bs{u}\times\bs{v}|}\) 时等号成立，故最大值为 \(\dfrac{4\sqrt3}{3}\)，选 B.
\end{solution}

\section{填空题}

\begin{problem}
不等式 \(\frac{x - 2}{x^2 + 3x + 2} > 0\) 的解集是\fillinblank{}.
\end{problem}

\begin{answer}
\(\left(-2,-1\right)\cup\left(2,+\infty\right)\)
\end{answer}

\begin{solution}
原不等式的定义域为 \(x\ne-2\)，\(x\ne-1\). 因式分解得
\[
\frac{x-2}{x^2+3x+2}=\frac{x-2}{(x+2)(x+1)}.
\]
在临界点 \(-2\)，\(-1\)，\(2\) 将数轴分成四个区间. 逐段判断分式符号可知，分式在 \((-2,-1)\) 和 \((2,+\infty)\) 上为正，在其余区间上为负. 因而解集为
\(\left(-2,-1\right)\cup\left(2,+\infty\right)\).
\end{solution}

\begin{problem}
已知 \(\alpha\) 为第二象限的角，\(\sin \alpha = \frac{3}{5}\)，则 \(\tan 2\alpha =\)\fillinblank{}.
\end{problem}

\begin{answer}
\(-\frac{24}{7}\)
\end{answer}

\begin{solution}
因为 \(\alpha\) 是第二象限角，所以
\(\cos\alpha=-\sqrt{1-\sin^2\alpha}=-\frac45\)，从而
\(\tan\alpha=\dfrac{\sin\alpha}{\cos\alpha}=-\dfrac34\).
于是
\[
\tan2\alpha=\frac{2\tan\alpha}{1-\tan^2\alpha}
=\frac{-\frac32}{1-\frac9{16}}
=-\frac{24}{7}.
\]
\end{solution}

\begin{problem}
某学校开设 \(A\) 类选修课 \(3\text{ 门}\)，\(B\) 类选修课 \(4\text{ 门}\)，一位同学从中共选 \(3\text{ 门}\)，若要求两类课程中各至少选一门，则不同的选法共有\fillinblank{}\(\text{种}\).
\end{problem}

\begin{answer}
30
\end{answer}

\begin{solution}
从 \(7\) 门课程中任取 \(3\) 门，再去掉全部选 \(A\) 类或全部选 \(B\) 类的情形，得到
\[
\mathrm C_7^3-\mathrm C_3^3-\mathrm C_4^3=35-1-4=30.
\]
也可按两类课程分别选取 \(1\) 门和 \(2\) 门，或 \(2\) 门和 \(1\) 门计数，结果相同. 所以不同的选法共有 \(30\) 种.
\end{solution}

\begin{problem}
已知 \(F\) 是椭圆 \(C\) 的一个焦点，\(B\) 是短轴的一个端点，线段 \(BF\) 的延长线交 \(C\) 于点 \(D\)，且 \( \overrightarrow{BF} = 2\overrightarrow{FD} \)，则 \(C\) 的离心率为\fillinblank{}.
\end{problem}

\begin{answer}
\(\frac{\sqrt{3}}{3}\)
\end{answer}

\begin{solution}
设椭圆的长半轴为 \(a\)，短半轴为 \(b\)，焦距参数为 \(c\)，并取一个焦点
\(F=(c,0)\)，短轴端点 \(B=(0,b)\). 由
\(\overrightarrow{BF}=2\overrightarrow{FD}\) 得
\[
D=F+\frac12(F-B)=\left(\frac{3c}{2},-\frac b2\right).
\]
因为 \(D\) 在椭圆上，所以
\[
\frac{\left(\frac{3c}{2}\right)^2}{a^2}
+\frac{\left(-\frac b2\right)^2}{b^2}=1,
\]
即 \(\dfrac{9c^2}{4a^2}+\dfrac14=1\). 因而
\(\dfrac{c^2}{a^2}=\dfrac13\)，所以椭圆的离心率
\(e=\dfrac ca=\dfrac{\sqrt3}{3}\).
\end{solution}

\section{解答题}

\begin{problem}
记等差数列 \(\{a_{n}\}\) 的前 \(n\) 项的和为 \(S_{n}\)，设 \(S_{3} = 12\)，且 \(2a_{1}\)，\(a_{2}\)，\(a_{3} + 1\) 成等比数列，求 \(S_{n}\).
\end{problem}

\begin{solution}
设等差数列的首项为 \(a_1\)，公差为 \(d\). 由 \(S_3=12\) 得
\[
S_3=\frac32\left(2a_1+2d\right)=3(a_1+d)=12,
\]
所以 \(a_1+d=a_2=4\). 又
\(a_3=a_1+2d\)，且 \(d=4-a_1\)，从而 \(a_3=8-a_1\).

由 \(2a_1\)，\(a_2\)，\(a_3+1\) 成等比数列，得
\[
a_2^2=2a_1(a_3+1),
\]
即 \(16=2a_1(9-a_1)\)，整理得 \(a_1^2-9a_1+8=0\).
所以 \(a_1=1\) 或 \(a_1=8\).

当 \(a_1=1\) 时，\(d=3\)，于是
\[
S_n=\frac n2\left(2a_1+(n-1)d\right)
=\frac n2\left(2+3n-3\right)
=\frac12n(3n-1).
\]
当 \(a_1=8\) 时，\(d=-4\)，于是
\[
S_n=\frac n2\left(16-4(n-1)\right)=2n(5-n).
\]
故 \(S_n=\frac12n(3n-1)\) 或 \(S_n=2n(5-n)\).
\end{solution}

\begin{problem}
已知 \(\triangle ABC\) 的内角 \(A\)，\(B\) 及其对边 \(a\)，\(b\) 满足 \(a + b = a \cot A + b \cot B\)，求内角 \(C\).
\end{problem}

\begin{solution}
由正弦定理可设 \(a=2R\sin A\)，\(b=2R\sin B\)，其中 \(R\) 为三角形外接圆半径. 同理，
\(a\cot A=2R\cos A\)，\(b\cot B=2R\cos B\). 所以题设等式等价于
\[
\sin A+\sin B=\cos A+\cos B.
\]
利用和差化积公式，并注意 \(A+B\in(0,\pi)\)，有
\[
2\sin\frac{A+B}{2}\cos\frac{A-B}{2}
=2\cos\frac{A+B}{2}\cos\frac{A-B}{2}.
\]
由于 \(\left|\frac{A-B}{2}\right|<\frac{\pi}{2}\)，故
\(\cos\dfrac{A-B}{2}>0\)，可以约去，得到
\(\sin\dfrac{A+B}{2}=\cos\dfrac{A+B}{2}\). 又
\(\dfrac{A+B}{2}\in(0,\frac{\pi}{2})\)，所以
\(\dfrac{A+B}{2}=\dfrac{\pi}{4}\). 因而 \(A+B=\dfrac{\pi}{2}\)，从而
\(C=\pi-A-B=\dfrac{\pi}{2}\).
\end{solution}

\begin{problem}
投到某杂志的稿件，先由两位初审专家进行评审. 若能通过两位初审专家的评审，则予以录用；若两位初审专家都未予通过，则不予录用；若恰能通过一位初审专家的评审，则再由第三位专家进行复审，若能通过复审专家的评审，则予以录用，否则不予录用. 设稿件能通过各初审专家评审的概率均为 \(0.5\)，复审的稿件能通过评审的概率为 \(0.3\). 各专家独立评审.

\begin{enumerate}
\item 求投到该杂志的 \(1\text{ 篇}\)稿件被录用的概率.
\item 求投到该杂志的 \(4\text{ 篇}\)稿件中，至少有 \(2\text{ 篇}\)被录用的概率.
\end{enumerate}
\end{problem}

\begin{solution}
\begin{enumerate}
\item 两位初审专家都通过的概率为 \(0.5^2=0.25\). 恰有一位通过的概率为
\(2\times0.5\times0.5=0.5\)，这时还需复审并通过，概率为 \(0.5\times0.3=0.15\). 因而一篇稿件被录用的概率为
\[
p=0.25+0.15=0.40.
\]
\item 设 \(X\) 表示 \(4\) 篇稿件中被录用的篇数. 各篇稿件的评审相互独立，所以
\(X\sim B(4,0.4)\). 因此
\[
P(X\geqslant2)=1-P(X=0)-P(X=1)
=1-0.6^4-\mathrm C_4^1(0.4)(0.6)^3
=1-0.1296-0.3456=0.5248.
\]
\end{enumerate}
\end{solution}

\begin{problem}
如图，四棱锥 \(S - ABCD\) 中，\(SD \perp\) 底面 \(ABCD\)，\(AB \parallel DC\)，\(AD \perp DC\)，\(AB = AD = 1\)，\(DC = SD = 2\)，\(E\) 为棱 \(SB\) 上的一点，平面 \(EDC \perp\) 平面 \(SBC\).

\begin{enumerate}
\item 证明：\(SE = 2EB\).
\item 求二面角 \(A - DE - C\) 的大小.
\end{enumerate}

\bitmapfigure[max width=0.38\linewidth]{img/2010/syllabus_paper_1_liberal/q20_fig1.png}
\end{problem}

\begin{solution}
\begin{enumerate}
\item 建立空间直角坐标系，取 \(D\) 为原点，\(DC\) 沿 \(x\) 轴，\(DA\) 沿 \(y\) 轴，\(DS\) 沿 \(z\) 轴，则 \(D=(0,0,0)\)，\(C=(2,0,0)\)，\(A=(0,1,0)\)，\(B=(1,1,0)\)，\(S=(0,0,2)\).

设 \(t=\dfrac{SE}{SB}\)，则 \(0<t<1\)，且
\[
E=S+t(B-S)=(t,t,2-2t).
\]
平面 \(EDC\) 的两个方向向量可取
\(\overrightarrow{DC}=(2,0,0)\)，\(\overrightarrow{DE}=(t,t,2-2t)\)，所以其法向量可取
\(\bs{n}_1=(0,-2(1-t),t)\). 平面 \(SBC\) 的两个方向向量可取
\(\overrightarrow{SB}=(1,1,-2)\)，\(\overrightarrow{SC}=(2,0,-2)\)，其法向量可取
\(\bs{n}_2=(1,1,1)\).

由两平面垂直知 \(\bs{n}_1\cdot\bs{n}_2=0\)，即
\(-2(1-t)+t=0\). 解得 \(t=\dfrac23\)，从而
\(\frac{SE}{SB}=\frac23\)，\(\frac{EB}{SB}=1-\frac23=\frac13\)，
所以 \(SE=2EB\).
\item 由 \(t=\dfrac23\) 得 \(E=(\frac23,\frac23,\frac23)\)，故直线 \(DE\) 的方向向量可取
\(\bs{d}=(1,1,1)\). 在垂直于 \(DE\) 的截面内，二面角 \(A-DE-C\) 等于向量
\(\overrightarrow{DA}\)，\(\overrightarrow{DC}\) 在该截面上的投影所成的角.

将 \(\overrightarrow{DA}=(0,1,0)\)，\(\overrightarrow{DC}=(2,0,0)\) 分别投影到 \(\bs d^\perp\)，得
\[
\bs a=(0,1,0)-\frac{(0,1,0)\cdot\bs d}{|\bs d|^2}\bs d
=\left(-\frac13,\frac23,-\frac13\right),
\]
\[
\bs c=(2,0,0)-\frac{(2,0,0)\cdot\bs d}{|\bs d|^2}\bs d
=\left(\frac43,-\frac23,-\frac23\right).
\]
于是 \(\bs a\cdot\bs c=-\frac23\)，\(|\bs a|=\frac{\sqrt6}{3}\)，\(|\bs c|=\frac{2\sqrt6}{3}\).
所以
\[
\cos\angle A-DE-C=\frac{\bs a\cdot\bs c}{|\bs a||\bs c|}=-\frac12,
\]
故二面角 \(A-DE-C\) 的大小为 \(120^{\circ}\).
\end{enumerate}
\end{solution}

\begin{problem}
已知函数 \( f(x) = 3ax^4 - 2(3a + 1)x^2 + 4x \).

\begin{enumerate}
\item 当 \(a = \frac{1}{6}\) 时，求 \(f(x)\) 的极值.
\item 若 \(f(x)\) 在 \(\left(-1,1\right)\) 上是增函数，求 \(a\) 的取值范围.
\end{enumerate}
\end{problem}

\begin{solution}
\begin{enumerate}
\item 当 \(a=\dfrac16\) 时，
\[
f(x)=\frac12x^4-3x^2+4x.
\]
\[
f'(x)=2x^3-6x+4=2(x-1)^2(x+2).
\]
当 \(x<-2\) 时 \(f'(x)<0\)，当 \(x>-2\) 且 \(x\ne1\) 时 \(f'(x)>0\). 因此 \(f(x)\) 在 \(\left(-\infty,-2\right)\) 上递减，在 \((-2,1)\) 和 \((1,+\infty)\) 上递增；\(x=1\) 处导数虽为零，但导数不变号，不是极值点. 于是 \(x=-2\) 为极小值点，且
\[
f(-2)=\frac12\cdot16-3\cdot4+4(-2)=-12.
\]
所以 \(f(x)\) 的极小值为 \(-12\)，无极大值.
\item 对任意 \(a\)，有
\[
f'(x)=12ax^3-4(3a+1)x+4
=4(1-x)\bigl[1-3ax(x+1)\bigr].
\]
在 \((-1,1)\) 内 \(1-x>0\)，故 \(f(x)\) 在该区间上为增函数的充要条件是
\[
1-3ax(x+1)\geqslant0\quad(-1<x<1).
\]
函数 \(x(x+1)\) 在 \((-1,1)\) 上的最小值为 \(-\dfrac14\)，其上界为 \(2\).

当 \(a\geqslant0\) 时，最严格的限制来自 \(x\) 趋近 \(1\)，故需 \(1-6a\geqslant0\)，即 \(a\leqslant\dfrac16\). 当 \(a<0\) 时，令 \(a=-b\)，其中 \(b>0\)，则
\[
1-3ax(x+1)=1+3b x(x+1)\geqslant1-\frac{3b}{4}.
\]
所以需 \(b\leqslant\dfrac43\)，即 \(a\geqslant-\dfrac43\). 两种情形合并得
\[
a\in\left[-\frac43,\frac16\right].
\]
在两个端点处导数只在区间内个别点为零而不变号，函数仍为增函数，故端点可以取到.
\end{enumerate}
\end{solution}

\section{选考题：解答题}

\begin{problem}
已知抛物线 \(C: y^{2} = 4x\) 的焦点为 \(F\)，过点 \(K(-1,0)\) 的直线 \(l\) 与 \(C\) 相交于 \(A\)，\(B\) 两点，点 \(A\) 关于 \(x\) 轴的对称点为 \(D\).

\begin{enumerate}
\item 证明：点 \(F\) 在直线 \(BD\) 上.
\item 设 \(\overrightarrow{FA} \cdot \overrightarrow{FB} = \frac{8}{9}\)，求 \(\triangle BDK\) 的内切圆 \(M\) 的方程.
\end{enumerate}
\end{problem}

\begin{solution}
\begin{enumerate}
\item 抛物线 \(y^2=4x\) 上的点可参数表示为 \(P(t)=(t^2,2t)\). 设
\(A=(u^2,2u)\)，\(B=(v^2,2v)\)，其中 \(u\ne v\). 弦 \(AB\) 的斜率为
\[
\frac{2v-2u}{v^2-u^2}=\frac{2}{u+v},
\]
其方程为
\[
y=\frac{2}{u+v}x+\frac{2uv}{u+v}.
\]
因为直线 \(AB\) 经过 \(K(-1,0)\)，代入得
\[
0=-\frac{2}{u+v}+\frac{2uv}{u+v},
\]
所以 \(uv=1\).

点 \(A\) 关于 \(x\) 轴的对称点为 \(D=(u^2,-2u)\)，而焦点为 \(F=(1,0)\). 直线 \(BD\) 的斜率为
\[
\frac{2v+2u}{v^2-u^2}=\frac{2}{v-u}.
\]
又因为 \(v=\dfrac1u\)，所以
\[
\frac{2}{v-u}=\frac{2u}{1-u^2}
=\frac{0-(-2u)}{1-u^2},
\]
这正是直线 \(DF\) 的斜率. 因而 \(F\) 在直线 \(BD\) 上.
\item 由 \(uv=1\)，有
\(\overrightarrow{FA}=(u^2-1,2u)\)，\(\overrightarrow{FB}=(v^2-1,2v)\).
因此
\[
\overrightarrow{FA}\cdot\overrightarrow{FB}
=(u^2-1)(v^2-1)+4uv
=6-u^2-v^2.
\]
由题设点积为 \(\dfrac89\)，得 \(u^2+v^2=\frac{46}{9}\)，从而 \((u+v)^2=u^2+v^2+2uv=\frac{64}{9}\).
记 \(w=u+v\)，则 \(|w|=\dfrac83\). 上下对称的两种情形所得内切圆相同，以下取 \(w=\dfrac83\)，于是 \(u,v>0\).

由 \(K=(-1,0)\)，\(B=(v^2,2v)\)，\(D=(u^2,-2u)\)，并利用 \(uv=1\)，得
\[
KB^2=(v^2+1)^2+4v^2
=v^2\left(u^2+v^2+6\right)
=\frac{100v^2}{9},
\]
\[
KD^2=(u^2+1)^2+4u^2=\frac{100u^2}{9}.
\]
所以 \(KB=\dfrac{10v}{3}\)，\(KD=\dfrac{10u}{3}\). 又
\[
BD^2=(v^2-u^2)^2+4(u+v)^2
=(v-u)^2w^2+4w^2=w^4,
\]
故 \(BD=w^2=\dfrac{64}{9}\). 从而
\[
KB+KD+BD=\frac{10(u+v)}3+\frac{64}{9}
=\frac{80}{9}+\frac{64}{9}=16.
\]
三角形 \(BDK\) 的内心为
\[
I=\frac{KD\cdot B+KB\cdot D+BD\cdot K}{KD+KB+BD}.
\]
计算其坐标，得
\[
x_I=\frac{\frac{10u}{3}v^2+\frac{10v}{3}u^2-\frac{64}{9}}{16}
=\frac{\frac{10}{3}uv(u+v)-\frac{64}{9}}{16}
=\frac19,
\]
\[
y_I=\frac{\frac{10u}{3}(2v)+\frac{10v}{3}(-2u)}{16}=0.
\]
所以 \(I=(\frac19,0)\). 直线 \(BK\) 即直线 \(l\)，方程为
\[
2x-wy+2=0.
\]
点 \(I\) 到该直线的距离为
\[
r=\frac{\left|2\cdot\frac19+2\right|}{\sqrt{4+w^2}}
=\frac{\frac{20}{9}}{\frac{10}{3}}=\frac23.
\]
这是内心到三边的距离，即内切圆半径. 因而内切圆 \(M\) 的方程为
\[
\left(x-\frac19\right)^2+y^2=\frac49.
\]
\end{enumerate}
\end{solution}
